% Preambolo comune delle simulazioni d'esame (vedi GUIDA.md).
% Ogni esercizio: \input{../preambolo} e \graphicspath verso le figure dello
% svolgimento di riferimento; si compila con compila.sh (svolgimento.pdf, A4).

\documentclass[11pt]{article}

\usepackage[T1]{fontenc}
\usepackage[utf8]{inputenc}
\usepackage[italian]{babel}
\usepackage{amsmath,amssymb}
\usepackage{array}
\usepackage{siunitx}
\usepackage{booktabs}
\usepackage{graphicx}
\usepackage{float}
\usepackage{caption}
\usepackage{enumitem}
\usepackage{hyperref}
\usepackage[a4paper,margin=2.5cm]{geometry}

\hypersetup{hidelinks}
\setlength{\parindent}{0pt}
\setlength{\parskip}{4pt}
\allowdisplaybreaks
\sisetup{output-decimal-marker={,}, group-separator={\,}, per-mode=symbol,
         exponent-product=\cdot, range-phrase={\text{--}}}

% Notazione del formulario (come template/svolgimento.tex).
\newcommand{\Cp}{C^{+}}
\newcommand{\Cm}{C^{-}}
\newcommand{\Mr}{M_{r}}
\newcommand{\Ms}{M_{s}}

% Separatore fra relazioni nella stessa formula: \begin{gather*} A \ee B \end{gather*}.
\newcommand{\ee}{,\qquad}

% Titolo dell'esercizio: \esercizio{Prova del GG/MM/AAAA -- Esercizio n}{argomento}.
\newcommand{\esercizio}[2]{%
  \title{#1\\[2pt]\large #2}\author{}\date{}%
  \hypersetup{pdftitle={#1: #2}}}

% Figura centrata: \figura[frazione della larghezza]{file}{didascalia\label{...}}.
\newcommand{\figura}[3][0.7]{%
  \begin{figure}[H]\centering\includegraphics[width=#1\textwidth]{#2}\caption{#3}\end{figure}}

% Testo a sinistra e figura (anche un dettaglio ingrandito) accanto, per tenere la figura vicino al punto che
% si discute: \accanto[larghezza della figura]{file}{didascalia\label{...}}{testo}.
\newcommand{\accanto}[4][0.40]{%
  \par\noindent
  \begin{minipage}[t]{\dimexpr\textwidth-#1\textwidth-0.03\textwidth\relax}\vspace{0pt}#4\end{minipage}\hfill
  \begin{minipage}[t]{#1\textwidth}\vspace{0pt}\centering\includegraphics[width=\linewidth]{#2}%
  \captionof{figure}{#3}\end{minipage}\par\medskip}
